\documentclass[a4paper,12pt,hidelinks]{report}

%%Pacchetti utili anche se non necessari

\usepackage{amsfonts}
\usepackage{amsmath}
\usepackage{latexsym}
\usepackage{tabularx}
\usepackage[italian]{babel}
\usepackage[bookmarks=true]{hyperref}
\usepackage{url}
% \usepackage{subfigure}
\usepackage{epstopdf}
\usepackage[utf8]{inputenc}
% \usepackage[utf8x]{inputenc}
\usepackage{listings}
\usepackage{graphicx}
%-------------------------------------------

% \title{Progettazione sito web\\ ''B\&B La Vecchia Posta''}
% \author{Daniele Di Pompeo \\mat. 226766}
% \annoaccademico{2013-2014}
\begin{document}
  \begin{titlepage}
    \begin{center}
    % Upper part of the page
      \includegraphics[width=0.5\textwidth,keepaspectratio=true]{../img/logo}\\[1cm]    
      \textsc{\LARGE Covalide prototipo di navigazione}\\[0.6cm]
      \textsc{\LARGE  progetto del sito:\\[0.5cm] ``B\&B La Vecchia Posta''}\\ [2.0cm]

    % Author and supervisor
      \begin{minipage}{0.8\textwidth}
	\begin{flushleft} \large
	  \emph{Autore:} Daniele Di Pompeo \\[0.5cm]
	  \emph{Versione documento: 1.0}\\[0.5cm]
	  \emph{Data emissione del documento: \today}\\[0.5cm]
	\end{flushleft}
      \end{minipage}
    \end{center}
  \end{titlepage}

% \tableofcontents
 
\begin{abstract}
 In questo documento vengono elencate le peculiarità grafiche del sito web in oggetto che sono state descritte nel documento utente allegato alla documentazione 
 rilasciata al committente.
\end{abstract}

\section*{Layout generale}
Il protipo del layout utilizzato per la realizzazione del sito web è stato sottoposto al test di usabilità di vari tipi di utenti.
\begin{itemize}
 \item \textbf{Utente esperto}: l'utente esperto ci ha fornito un feedback positivo sia per quanto riguarda lo stile del sito sia per la facile navigazione tra le sezione
 dello stesso. In particolare è riuscito subito a raggiungere la sezione desiderata senza impiegare troppo tempo.
 \item \textbf{Navigatore esperto}: facendo utilizzare il prototipo del sito web ad un navigatore esperto (utente che trascorre alcune ore giornaliere su internet), il sito 
 risulta essere abbastanza navigabile e di gradevole impatto visivo.
 \item \textbf{Navigatore inesperto}: forse il test più importante per assicurare la buona riuscita della progettazione del sito. I risultati ottenuti si possono considerare 
 accetabili in relazione al nostro utente. 
 \par Nonostante la presenza di effetti grafici anche complessi il sito è risultato navigabile nella sua quasi totalità. L'utente 
 ha avuto difficoltà nell'utilizzare il form di prenotazione nell'homepage, non intuendo subito la possibilità di scelta della tipologia della camera. Probabilmente causata
 da una scelta non ottimale dell'elemento grafico. Si provvederà alla sostituzione dello stesso con uno più consono.
\end{itemize}

\section*{Struttura dei munu}
La struttura del menu risulta essere completamente visibile sia alla risoluzione scelta per la progettazione, sia ad una risoluzione più bassa.
Sono stati evitati scroll orizzontali nelle risoluzioni 1024px e 1200px (risoluzione di base) che coprono quasi la totalità delle configurazioni degli utenti che navigano
su internet.
Per quanto riguarda risoluzioni più basse e dispositivi mobile, si rimanda ad una versione successiva dello sviluppo.

\section*{Navigazione}
L'intera struttura delle pagine del sito maniene sempre la stessa struttura, come riportato nella progettazione delle gabbie logiche.
Le varie sezioni delle pagine sono state strutturate per ottenere una similitudine tra le varie sezioni. Si è utilizzata una prima pagina di resoconto delle varie sottosezioni
che permette l'individuazione dei vari contenuti. Nelle sezioni intermedie, la singola camera, il singolo evento, ecc. si è scelto di inserire uno slider relativo alla sezione
o di elimare lo slider ed inserire il widget della mappa, fornendo a destra del widget sia il menu di secondo livello che una scorciatoia alla prenotazione. Appena sopra la piega
appare il titolo della sezione per invogliare l'utente a visionare il resto della pagina. Come breadcrump si è scelto di evidenziare solamente la macro sezione in cui ci si trova
e non di mostrare una breadcrump visiva da inserire al di sotto del menu, cosa difficile da inserire causa effetto a tendina verticale del menu.

\section*{Etichette}
Le etichette del menu rispettano quasi 1a1 le sezioni individuate nel documento dei requisiti (capitolo x sezione y) e sono risultate comprensibili agli utenti che hanno eseguito
il primo test di usabilità.
Le etichette inoltre risultano essere coerenti con i contenuti delle sezioni che identificano. 
Ciò è testimoniato anche dai vari test effettuati.

\end{document}          
